\chapter{汇率与美元体系：为什么全世界用别人的钱做生意}
\chaptermeta{12}

\begin{ledetext}
一个宁波老板卖货给巴西人，两边都不是美国人，合同却用美元签。这背后是一套没人投票选出来的秩序。
\end{ledetext}

\section{一块钱值多少钱，得问另一块钱}

老韩 2011 年从一家台资厂出来单干，在宁波北仑做汽车密封条，厂里现在 63 个人。2024 年他接了圣保罗一家改装车厂的单子，132 万美元，账期 90 天。

报价那天美元兑人民币 7.24，钱到账那天 7.08。同样 132 万美元，少收了 21 万人民币。

他那个月的利润，就在这 16 个基点上没了。

汇率不是一个国家的考试成绩，它是两种钱互相的报价。"7.1"的意思是 1 美元换 7.1 元人民币，同一个数倒过来读，就是 1 元人民币换 0.141 美元。两个方向，谁也不比谁高级。

谁在决定这个数？

老韩的第一反应是贸易：中国出口多，外国人得先买人民币才能付货款，人民币就该涨。这个逻辑本身没错，可它解释不了他看到的行情。

全球外汇市场日均成交是万亿美元量级。而全球货物贸易一整年的总额，摊到每个工作日也就几百亿美元的量级。

\textbf{在外汇市场里，做真实生意的钱是少数派。}

绝大部分成交来自资本流动：跨境买卖股票债券、跨境放贷、对冲敞口、套利、投机。它们盘子大、掉头快、几乎没有摩擦成本。老韩那 132 万美元排队等着结汇的时候，隔壁有人几秒钟就调走了十亿。

所以看汇率，先看钱往哪儿跑，再看货往哪儿走。一个国家完全可能贸易顺差很大而货币在贬，也完全可能常年逆差而货币很硬。美国已经这么硬了几十年。

\begin{chainflow}
\chainnode{美国加息}\chainarrow \chainnode{美元资产更香}\chainarrow \chainnode{资本流入美国}\chainarrow \chainnode{美元走强}\chainarrow \chainnode{新兴市场货币承压}\chainarrow \chainnode{他们的美元债还起来更贵}
\end{chainflow}

\section{巨无霸能告诉你方向，不能告诉你时间}

《经济学人》1986 年发明了巨无霸指数，逻辑简单到可爱：同一个汉堡，在哪儿都该值一样多的钱。

假设一个巨无霸在美国卖 5.8 美元，在中国卖 25 元，"汉堡汇率"就是 25 ÷ 5.8 ≈ 4.3。要是市场汇率在 7 以上，按汉堡算，人民币被低估了四成。

这是\term{购买力平价}{purchasing power parity}\index{021@购买力平价}的直觉：长期看，汇率该让同一篮子东西在各国卖出差不多的价。

麻烦在于，汉堡运不过去。

一个巨无霸里，牛肉和面包是可以跨境运的，可门店租金、店员工资、水电、当地税全都运不走。上海店员的工资和芝加哥店员的工资差多少，汉堡价格就差多少，套利者没办法把上海的店员打包空运到芝加哥去。

这个漏洞可以一直漏下去。加上关税、增值税、饮食习惯，PPP 算出来的偏差能常年存在，一存几十年。

它是一根很长的橡皮筋，不是一张时刻表。谁拿它判断"下个季度该升还是该贬"，等于拿望远镜看牙齿。

\section{借日元，买美债：全世界最拥挤的那笔生意}

\begin{egbox}{举个栗子}
2026 年 8 月，同一个地球上，钱的价格差出好几倍。日本央行政策利率 \textbf{1.0\%}，2026 年 6 月刚从 0.75\% 加上来，是 1995 年 9 月以来最高；美联储联邦基金利率目标区间 \textbf{3.50\%–3.75\%}；中国 1 年期 LPR \textbf{3.0\%}。

于是有人做这样一笔生意：在东京借日元，成本大约 1\%；换成美元，买美国短期国债，收益大约 3.6\%。什么活都不用干，一年白赚 2.6 个百分点。

2.6\% 不解渴，那就上杠杆。5 倍，变 13\%。10 倍，变 26\%。

\end{egbox}

这就是\term{套息交易}{carry trade}\index{058@套息交易}，全球资金流动里最庞大也最不起眼的一股力量。它的规模没人能精确统计，因为它藏在无数张互不相干的资产负债表里。

它的毛病只有一句话：\textbf{利差是按天赚的，汇率是按小时还的。}

\begin{egbox}{举个栗子：2024 年 8 月那三天}
2024 年 7 月 31 日，日本央行把政策利率从 0～0.1\% 上调到 0.25\%，同时缩减购债。加息 25 个基点，听起来像挠痒痒。

但成本端动了。日元开始升值，美元兑日元从 7 月初的 161 一路走到 8 月初的 142 附近，不到一个月，日元强了十几个百分点。

借日元的人要还钱，得用更贵的日元买回来。攒了好几年的那点利差，被汇率一口吞掉。

平仓引发平仓。8 月 5 日，日经 225 单日跌 12.4\%，创下它历史上最大的单日点数跌幅；美股同步大跌，波动率指数盘中冲到 60 以上。从东京传到纽约，用了几个小时。

起因是日本一次 25 个基点的加息。

\end{egbox}

\begin{pullquote}{}
套息交易的收益曲线像爬楼梯，风险曲线像坐电梯。三年爬上去，三天掉下来。
\end{pullquote}

\section{三选二：所有国家都在这张桌子上做同一道题}

\term{三元悖论}{impossible trinity}\index{045@三元悖论}说的是：固定汇率、资本自由流动、独立货币政策，这三样你最多同时拿两样。

为什么？假设你既要汇率钉死，又要资本自由进出。美联储把利率抬到 3.75\%，你的利率是 1\%，资本立刻往美国跑，你的货币面临贬值压力。可你承诺了固定汇率，就只剩两条路：拿外汇储备去买回本币，或者把自己的利率也抬到 3.75\%。前者会把储备耗光，后者等于承认你的货币政策是美国定的。

三样都想要的，历史上没有一个撑住。1992 年英镑被迫退出欧洲汇率机制，1997 年泰铢崩盘，做错的都是同一道题。

\begin{tablewrap}
\begin{xltabular}{\linewidth}{>{\raggedright\arraybackslash}p{0.20\linewidth}XXX}
\toprule
\bthead{经济体} & \bthead{放弃了什么} & \bthead{换来了什么} & \bthead{代价长什么样} \\
\midrule
\textbf{香港} & 独立货币政策 & 港元钉住美元在 7.75–7.85 区间，资本完全自由进出 & 美联储加息，香港必须跟着加，哪怕本地楼市和零售正在挨打 \\
\textbf{欧元区成员国} & 独立货币政策 & 成员国之间等于彻底固定汇率（用同一种钱），资本自由流动 & 2010–2012 年主权债务危机时，希腊没法靠贬值恢复竞争力，只能靠内部通缩，也就是直接砍工资 \\
\textbf{美国、日本} & 固定汇率 & 完全独立的货币政策，资本自由流动 & 汇率大幅波动。2024 年 8 月日元三周走强十几个百分点，全球市场跟着地震 \\
\textbf{中国} & 资本完全自由流动 & 独立的货币政策，有管理的浮动汇率 & 跨境资金进出有额度和审批，老韩换汇要提交合同和报关单 \\
\bottomrule
\end{xltabular}
\end{tablewrap}

香港金管局跟着美联储走不是没主见，是合同条款；欧洲央行一套利率管二十个国家，必然有人被夹疼，这是设计使然；中国能在美联储加息的时候降息，靠的就是老韩填的那些表格。

\section{老韩为什么不用人民币报价}

他试过。巴西客户第一反应是"那我的银行怎么给我报价"。

\begin{flowlist}
  \flowitem{01}{网络效应}{老韩的橡胶原料按美元报价，海运费按美元结算，货运险和信用证也是美元。他一个人换币种，等于把自己从整张网里摘出来。}
  \flowitem{02}{有个足够深的池子停钱}{收到美元不能压箱底，得有地方放。美国国债可流通市场规模超过 30 万亿美元，今天卖 10 亿也砸不动价格。这种深度全世界只有这一家。}
  \flowitem{03}{随时能换出来}{这一条比"收益高"重要得多。储备货币的第一要求是想走的时候走得掉。}
  \flowitem{04}{法治与合约的可预期性}{钱放在纽约，出了纠纷有一套大致能预判的司法程序。这条听着虚，定价的时候非常实。}
  \flowitem{05}{清算管道}{SWIFT 报文加上美元清算体系，是跨境支付的默认路径。用了半个世纪，改造成本高到没人愿意先动。}
\end{flowlist}

\begin{deepbox}{进阶：可以跳过}
1960 年，比利时经济学家罗伯特·特里芬指出了一个绕不过去的矛盾。

全世界都想用美元做储备和结算，那美元就必须源源不断流到美国之外去。怎么送出去？靠贸易逆差。美国进口比出口多，把美元付给全世界。

可长期逆差意味着美国持续对外欠债。欠得越多，别人越怀疑这些美元将来还值不值这个价。

\textbf{世界需要美元的方式，恰恰会削弱美元。}这就是特里芬难题，它没有解，只有拖。布雷顿森林体系 1971 年就是被这道题拖垮的，当时的表现形式是美国的黄金储备不够兑付境外的美元。今天它换了张脸：美国国债在 2026 年 8 月突破 40 万亿美元，IIF 警告这条路"日益走上不可持续的道路"。

\end{deepbox}

\section{2026：储备资产的第一属性变了}

欧洲央行 2026 年 6 月的《欧元的国际角色》报告，给了一组数字。

\begin{barfig}
  \barrow{黄金}{27}{27\%}
  \barrow{美国国债}{22}{22\%}
  \barrow{其他美元资产}{20}{20\%}
  \barrow{欧元}{15}{15\%}
\end{barfig}

截至 2025 年底，黄金占全球央行储备的 27\%，美国国债占 22\%。\textbf{黄金第一次超过美债，成为各国央行的第一大储备资产。}2024 年底这个对比还是黄金 20\%、美债 25\%，其他美元资产从 22\% 降到 20\%，欧元稳在 15\%。

美元计价资产合计仍占 42\%，依然是最大的单一币种阵营。这不是崩塌，是配比在挪。

\begin{statrow}{3}
  \statitem{3.6}{万吨+}{全球央行黄金储备，接近布雷顿森林体系鼎盛期的约 3.8 万吨}
  \statitem{850}{吨}{2025 年全球央行净购金}
  \statitem{5600}{美元/盎司}{2026 年 1 月金价一度接近的水平}
\end{statrow}

为什么突然开始抢金子？

分水岭是 2022 年。俄乌冲突之后，西方国家冻结了俄罗斯央行数千亿美元规模的外汇储备。

在那之前，各国央行管钱的考核表上写着三项：安全、流动性、收益。"会不会被没收"这一项根本不在表上。

那天之后，所有非西方阵营的储备管理者都做了同一道算术：我持有的美债，本质上是别人账本上的一行数字。那行数字能不能被划掉，不取决于我。

黄金没有这个问题。它是唯一一种不构成任何人负债的储备资产。搬进自己的金库，它就在那儿，不需要谁替你记账。代价是它不付利息，还得花钱看守，一吨金子一年的保管费不是小数目。

\begin{pullquote}{}
储备资产的第一属性，从"收益"变回了"不会被没收"。
\end{pullquote}

\begin{databox}{2026 数据：一个有点荒诞的细节}
2025 年全球最大的单一黄金买家不是哪国央行，是稳定币发行商 Tether，一年买进超过 100 吨。

一家专门给美元稳定币做储备的公司，拿赚来的钱去买金条。连最"美元原生"的那批玩家，也在给自己的储备垫一层不依赖任何账本的东西。

\end{databox}

\section{人民币走到哪一步了}

跨境贸易结算里人民币的占比这些年在稳步上升，CIPS 的参与机构和覆盖国家持续扩大，数字人民币的跨境支付网络 CBETS 在 2026 年新增了 26 家银行。老韩现在有两个南美客户愿意用人民币结算，三年前一个都没有。

这些是实打实的进展。但真正的约束不在管道上。

\textbf{第一，资本项目还没有完全可兑换。}别人愿意拿一种货币当储备的前提，是想走的时候能走。这条和三元悖论直接顶上：中国选了独立货币政策加有管理的汇率，就必须留着资本管制这道闸。要放开，就得在另外两样里让出一样。这是个真实的取舍，不是意愿问题。

\textbf{第二，国债市场的深度和开放度。}储备货币需要一个巨大的、外国人能自由进出的、买卖不怎么影响价格的安全资产池。中国国债市场的规模已经不小，但外资持有比例和交易便利程度，跟美债还不在一个量级。

什么时候能实现资本项目基本可兑换？这件事目前没人知道答案，包括正在做这个决定的人。它取决于国内金融体系的稳健程度、外部环境、以及那一年的风险偏好。

\begin{mythbox}{常见误解}
\mvfrow{误解}{"本币贬值就是坏事，升值说明国力强。"}
\mvfrow{事实}{得分三条线看。\textbf{出口企业：}人民币贬值，老韩那 132 万美元换回更多人民币，报价也更有竞争力，他是受益方。\textbf{进口方和你我：}贬值让原油、芯片、留学、出国旅游全部变贵，输入型通胀跟着来。\textbf{有外币负债的：}这条最凶险。老韩 2023 年买德国挤出机借了 260 万美元设备贷，但他七成收入是美元，这叫自然对冲，扛得住。扛不住的是那种只赚人民币却借了美元债的公司：人民币贬 10\%，收入一分没变，债务折成人民币凭空多出一大块。1997 年亚洲金融危机里倒下的那批企业，死因高度一致，行话叫\term{币种错配}{currency mismatch}\index{005@币种错配}。}
\mvfrow{误解}{"美元马上要崩溃了" / "人民币很快就会取代美元"。}
\mvfrow{事实}{这两句话都是错的，而且错在同一个地方：把趋势的方向当成了趋势的速度。截至 2025 年底，美元计价资产仍占全球央行储备的 42\%，是最大的单一阵营；同时黄金一年之内从 20\% 升到 27\%，2025 年欧元计价国际债券发行增长 30\%、接近 1 万亿欧元的历史新高，国际投资者净流入欧元区资产 8500 亿欧元。\hlmark{这叫渐进多元化，不叫替代。}储备货币的更替以几十年计：英镑让位给美元花了大半个世纪，中间还隔着两次世界大战。}
\end{mythbox}

本章不预测任何货币的走势，也不给任何换汇建议。老韩的锁汇决定是他自己做的，他做对过，也做错过。

\begin{takeaway}{本章一句话}
汇率是两种钱的相对价格，今天主要由资本流动而不是贸易决定；美元体系靠的是网络效应、市场深度和清算管道，而 2026 年正在发生的是渐进多元化。

\begin{itemize}
  \item 购买力平价是长期引力，不是短期指南。它告诉你方向，不告诉你时间。
  \item 套息交易赚的是利差，还的是汇率。2024 年 8 月日经单日跌 12.4\%，起因是日本加息 25 个基点。
  \item 三元悖论三选二：香港和欧元区成员国放弃独立货币政策，美日放弃固定汇率，中国放弃资本完全自由流动。
  \item 截至 2025 年底黄金占全球央行储备 27\%，首次超过美债的 22\%；美元计价资产合计仍占 42\%。
\end{itemize}

\end{takeaway}

\begin{bridgebox}
\textbf{下一章}到这里，管道全部铺好了：钱怎么造出来、怎么定价、怎么跨境流动。可管道里流的东西，为什么隔几年就要炸一次？下一层我们抬头看天气。
\end{bridgebox}

